% GENERATED FILE. Edit canonical lessons, then run npm run generate:books.
% canonical-source-hash: fnv1a64:f4a527469db64db8
% canonical-lessons: PA-C131-battakh, PA-C131-kirli, PA-C131-kacchu, PA-C131-magarmacch, PA-C131-macchar

\chapter{Duck, Lizard, Turtle, Crocodile, Mosquito}
\label{ch:pa-duck-lizard-turtle-crocodile-mosquito}
\hlchaptermodality{131}

\begin{chapteropening}
\textbf{By the end of this chapter:} \emph{I can say a duck, a lizard, a turtle, a crocodile and a mosquito.}

Complete the last lesson of chapter 131: i can say a duck, a lizard, a turtle, a crocodile and a mosquito.
\end{chapteropening}

\section[\texorpdfstring{\pa{ਬੱਤਖ਼} (battakh)}{ਬੱਤਖ਼ (battakh)}]{\pa{ਬੱਤਖ਼} (battakh) — a duck}

\label{lesson:PA-C131-battakh}



\begin{quote}
Before the new one: say the Punjabi for a wolf, then the Punjabi for a pigeon.
\end{quote}

\subsection*{You'll want to know: \pa{ਬੱਤਖ਼}}
\textbf{\pa{ਬੱਤਖ਼}} — \emph{battakh} — \textquotedblleft{}a duck\textquotedblright{}.

\begin{grammarlens}[title={What you've built}]
One more word for this chapter.
\end{grammarlens}

\subsection*{Your turn}
\begin{itemize}
\raggedright
  \item \emph{Say it:} \emph{battakh}
  \item \emph{Say it:} \emph{battakh}, once more
  \item \emph{Say it:} say \emph{battakh}
  \item \emph{Recall:} say the Punjabi for a wolf, then the Punjabi for a pigeon, then say \emph{battakh} again
\end{itemize}

\begin{tcolorbox}[breakable,colback=teal!4,colframe=teal!35!black,title={Before you move on}]
What does \pa{ਬੱਤਖ਼} mean? (\textquotedblleft{}A duck\textquotedblright{}.) Say it once more.
\end{tcolorbox}
\section[\texorpdfstring{\pa{ਕਿਰਲੀ} (kirlī)}{ਕਿਰਲੀ (kirlī)}]{\pa{ਕਿਰਲੀ} (kirlī) — a lizard}

\label{lesson:PA-C131-kirli}



\begin{quote}
Before the new one: say the Punjabi for a pigeon, then the Punjabi for a duck.
\end{quote}

\subsection*{You'll want to know: \pa{ਕਿਰਲੀ}}
\textbf{\pa{ਕਿਰਲੀ}} — \emph{kirlī} — \textquotedblleft{}a lizard\textquotedblright{}.

\begin{grammarlens}[title={What you've built}]
One more word for this chapter.
\end{grammarlens}

\subsection*{Your turn}
\begin{itemize}
\raggedright
  \item \emph{Say it:} \emph{kirlī}
  \item \emph{Say it:} \emph{kirlī}, once more
  \item \emph{Say it:} say \emph{kirlī}
  \item \emph{Recall:} say the Punjabi for a pigeon, then the Punjabi for a duck, then say \emph{kirlī} again
\end{itemize}

\begin{tcolorbox}[breakable,colback=teal!4,colframe=teal!35!black,title={Before you move on}]
What does \pa{ਕਿਰਲੀ} mean? (\textquotedblleft{}A lizard\textquotedblright{}.) Say it once more.
\end{tcolorbox}
\section[\texorpdfstring{\pa{ਕੱਛੂ} (kacchū)}{ਕੱਛੂ (kacchū)}]{\pa{ਕੱਛੂ} (kacchū) — a turtle}

\label{lesson:PA-C131-kacchu}



\begin{quote}
Before the new one: say the Punjabi for a duck, then the Punjabi for a lizard.
\end{quote}

\subsection*{You'll want to know: \pa{ਕੱਛੂ}}
\textbf{\pa{ਕੱਛੂ}} — \emph{kacchū} — \textquotedblleft{}a turtle\textquotedblright{}.

\begin{grammarlens}[title={What you've built}]
One more word for this chapter.
\end{grammarlens}

\subsection*{Your turn}
\begin{itemize}
\raggedright
  \item \emph{Say it:} \emph{kacchū}
  \item \emph{Say it:} \emph{kacchū}, once more
  \item \emph{Say it:} say \emph{kacchū}
  \item \emph{Recall:} say the Punjabi for a duck, then the Punjabi for a lizard, then say \emph{kacchū} again
\end{itemize}

\begin{tcolorbox}[breakable,colback=teal!4,colframe=teal!35!black,title={Before you move on}]
What does \pa{ਕੱਛੂ} mean? (\textquotedblleft{}A turtle\textquotedblright{}.) Say it once more.
\end{tcolorbox}
\section[\texorpdfstring{\pa{ਮਗਰਮੱਛ} (magarmacch)}{ਮਗਰਮੱਛ (magarmacch)}]{\pa{ਮਗਰਮੱਛ} (magarmacch) — a crocodile}

\label{lesson:PA-C131-magarmacch}



\begin{quote}
Before the new one: say the Punjabi for a lizard, then the Punjabi for a turtle.
\end{quote}

\subsection*{You'll want to know: \pa{ਮਗਰਮੱਛ}}
\textbf{\pa{ਮਗਰਮੱਛ}} — \emph{magarmacch} — \textquotedblleft{}a crocodile\textquotedblright{}.

\begin{grammarlens}[title={What you've built}]
One more word for this chapter.
\end{grammarlens}

\subsection*{Your turn}
\begin{itemize}
\raggedright
  \item \emph{Say it:} \emph{magarmacch}
  \item \emph{Say it:} \emph{magarmacch}, once more
  \item \emph{Say it:} say \emph{magarmacch}
  \item \emph{Recall:} say the Punjabi for a lizard, then the Punjabi for a turtle, then say \emph{magarmacch} again
\end{itemize}

\begin{tcolorbox}[breakable,colback=teal!4,colframe=teal!35!black,title={Before you move on}]
What does \pa{ਮਗਰਮੱਛ} mean? (\textquotedblleft{}A crocodile\textquotedblright{}.) Say it once more.
\end{tcolorbox}
\section[\texorpdfstring{\pa{ਮੱਛਰ} (macchar)}{ਮੱਛਰ (macchar)}]{\pa{ਮੱਛਰ} (macchar) — a mosquito}

\label{lesson:PA-C131-macchar}



\begin{quote}
Before the new one: say the Punjabi for a turtle, then the Punjabi for a crocodile.
\end{quote}

\subsection*{You'll want to know: \pa{ਮੱਛਰ}}
\textbf{\pa{ਮੱਛਰ}} — \emph{macchar} — \textquotedblleft{}a mosquito\textquotedblright{}.

\begin{grammarlens}[title={What you've built}]
One more word for this chapter.
\end{grammarlens}

\subsection*{Your turn}
\begin{itemize}
\raggedright
  \item \emph{Say it:} \emph{macchar}
  \item \emph{Say it:} \emph{macchar}, once more
  \item \emph{Say it:} say \emph{macchar}
  \item \emph{Recall:} say the Punjabi for a turtle, then the Punjabi for a crocodile, then say \emph{macchar} again
\end{itemize}

\begin{tcolorbox}[breakable,colback=teal!4,colframe=teal!35!black,title={Before you move on}]
What does \pa{ਮੱਛਰ} mean? (\textquotedblleft{}A mosquito\textquotedblright{}.) Say it once more.
\end{tcolorbox}
